\documentclass[10pt,conference]{IEEEtran}

\usepackage[utf8]{inputenc}
\usepackage[T1]{fontenc}
\usepackage{graphicx}
\usepackage{booktabs}
\usepackage{multirow}
\usepackage{amsmath}
\usepackage{amssymb}
\usepackage{hyperref}
\usepackage{url}
\usepackage{xcolor}
\usepackage{array}

\title{Kolmogorov-Arnold Networks for Forecasting:\\
An Interpretable Comparative Study on Climate and Electricity Data}

\author{
\IEEEauthorblockN{Kirill Alekhin}
\IEEEauthorblockA{
ITMO University\\
Saint Petersburg, Russia\\
\texttt{alekkhin.kirill.list.ru}
}
}

\begin{document}

\maketitle

\begin{abstract}
This paper presents an extended empirical study of Kolmogorov-Arnold Networks (KAN) and hybrid KAN architectures in forecasting tasks of different structure and complexity. Two case studies are considered. The first one addresses multivariate climate time series forecasting on the UCI Air Quality dataset, where KAN-based models are compared with recurrent and feed-forward neural baselines in both standard and robustness-oriented setups. The second one focuses on short-term electricity price forecasting on a leakage-safe panel dataset constructed from Russian regional energy market data. In the baseline electricity setup, tree-based models outperform KAN-based architectures by a large margin. However, after reformulating the forecasting target as a next-step price increment, introducing region embeddings, and applying residual learning, the strongest KAN-based model approaches the performance of the best boosting baseline. Across both case studies, the experiments show that the main value of KAN is not universal superiority in predictive accuracy, but the combination of nonlinear modeling flexibility and interpretability through learned channel-wise $\phi$-functions. The paper highlights both the strengths and the limits of KAN on real-world forecasting problems and provides a reproducible analysis framework for studying the trade-off between accuracy and interpretability.
\end{abstract}

\begin{IEEEkeywords}
Kolmogorov-Arnold Networks, forecasting, time series, panel data, electricity price, climate data, interpretability, hybrid neural models
\end{IEEEkeywords}

\section{Introduction}

Forecasting is one of the central tasks in applied machine learning and data mining. In practice, forecasting models are used to support operational planning in climate science, industrial systems, electricity markets, logistics, finance, and many other domains. The quality of forecasts affects decision making directly, which means that model errors can translate into operational, economic, or even safety-related risks. For this reason, modern forecasting research increasingly focuses not only on predictive accuracy, but also on stability, robustness, and interpretability.

Real-world forecasting tasks are difficult for at least three reasons. First, the observed dynamics are often nonlinear and partially nonstationary. Second, the available data may contain missing values, noise, duplicated signals, and heterogeneous feature types. Third, the forecasting objective itself can vary significantly: some tasks are essentially multivariate sequence modeling, while others are better seen as panel or tabular prediction with temporal dependencies. This diversity makes it unlikely that a single model family will dominate in all scenarios.

Recently, Kolmogorov-Arnold Networks (KAN) have emerged as a promising alternative to traditional multilayer perceptrons \cite{liu2024kan}. The key architectural idea of KAN is to learn one-dimensional nonlinear transformations on connections rather than relying only on fixed activations applied to hidden node outputs. This makes KAN interesting for forecasting for two reasons. On the one hand, it offers a flexible way to model nonlinear dependencies. On the other hand, it introduces a potentially interpretable representation through learned $\phi$-functions, which can be visualized and analyzed channel by channel.

Despite this promise, the role of KAN in practical forecasting remains insufficiently understood. Existing discussions often focus on benchmark or synthetic settings, while applied forecasting tasks involve substantially more constraints: leakage from future information, strong seasonal structure, panel heterogeneity, and the need to compare not only with neural baselines but also with strong non-neural models such as gradient boosting and linear regression.

This paper is motivated by the following research question: \emph{under what conditions can KAN and hybrid KAN architectures serve as effective and interpretable forecasting models on real data?} To answer this question, we investigate two forecasting problems with very different structure:
\begin{itemize}
    \item a multivariate climate time series forecasting task based on the UCI Air Quality dataset;
    \item a short-term electricity price forecasting task on Russian regional energy market data, formulated as a leakage-safe panel forecasting problem.
\end{itemize}

The first case study is useful because it provides a relatively clean multivariate sequence modeling setup, where KAN can be compared directly with LSTM and MLP models and where $\phi$-functions are easy to interpret. The second case study is substantially harder and more realistic: it combines hourly, daily, and monthly information, exhibits strong regional heterogeneity, and requires careful handling of leakage from future information.

The contributions of the paper are the following:
\begin{itemize}
    \item We provide a comparative KAN study on two applied forecasting domains with different structural properties.
    \item We construct a leakage-safe electricity forecasting dataset from Russian regional energy data and evaluate KAN-based and non-KAN baselines on it.
    \item We show that a naive KAN formulation is insufficient for the electricity task, but that target reformulation, region embeddings, and residual learning substantially improve KAN performance.
    \item We analyze KAN interpretability through $\phi$-functions, channel contributions, and grouped feature importance.
    \item We discuss the conditions under which KAN is most useful as a forecasting model and as an interpretable nonlinear analysis tool.
\end{itemize}

\section{Related Work}

\subsection{Kolmogorov-Arnold Representation and KAN}

The theoretical motivation of KAN stems from the Kolmogorov-Arnold representation theorem \cite{kolmogorov1957,arnold1963}, which states that multivariate continuous functions can be represented through superpositions of one-dimensional functions. This result has long been interpreted as a theoretical basis for constructing function approximation schemes with explicit one-dimensional nonlinear components.

Modern KAN architectures translate this principle into neural computation by replacing part of the traditional node-centric nonlinear processing with learned nonlinear edge transformations \cite{liu2024kan}. Instead of applying a fixed nonlinearity such as ReLU after a linear layer, KAN learns a family of basis-weighted one-dimensional functions that can approximate different forms of nonlinear input-output dependence.

\subsection{Forecasting Architectures}

For time series forecasting, recurrent and sequence models remain important baselines. LSTM networks are still widely used as strong recurrent baselines for multivariate temporal prediction, especially in moderate-sized settings where interpretability is not the primary goal. Transformer-based models have also become prominent in long-horizon forecasting \cite{zhou2021informer}, but in many applied scenarios simpler baselines remain difficult to outperform consistently.

In contrast, for tabular or panel forecasting tasks, tree-based ensembles and regularized linear models are often extremely competitive. Gradient boosting models are particularly strong when the data contain mixed-scale signals, nonlinear interactions, and moderate amounts of engineered lag and calendar features. Therefore, any applied claim about KAN must be evaluated not only against neural baselines such as LSTM or MLP, but also against strong tabular baselines.

The continued relevance of LSTM models in forecasting is well established in sequence modeling literature \cite{hochreiter1997lstm}. At the same time, the success of gradient boosting for tabular machine learning is equally well documented \cite{friedman2001greedy,prokhorenkova2018catboost}. This duality matters for our paper: the climate task is close to the type of problem where recurrent or neural sequence models are natural, whereas the electricity task contains enough engineered structure that strong tabular methods should be treated as first-class competitors rather than auxiliary baselines.

\subsection{Interpretability in Forecasting}

Interpretability is a major challenge in forecasting. Standard deep sequence models often act as black boxes, making it difficult to understand which variables, lags, or temporal regimes are driving predictions. One of the attractive aspects of KAN is that it naturally exposes nonlinear channel-wise transformations via learned $\phi$-functions. This does not make all KAN models fully transparent, but it provides a more structured interpretability mechanism than purely latent representations.

Interpretability is especially important in applied energy systems, where model outputs may be used for dispatch, cost planning, and market-oriented decision support. In such settings, a small gain in accuracy does not automatically dominate a model that is slightly less accurate but easier to audit, debug, and communicate to domain experts. This is one of the central practical motivations behind our comparative study.

\section{Methodology}

\subsection{Problem Formulation}

We consider supervised forecasting under two related but distinct formulations.

\textbf{Multivariate time series forecasting.} Given a sequence window
\begin{equation}
X_t = [x_{t-w+1}, \dots, x_t],
\end{equation}
where each $x_t \in \mathbb{R}^d$, the task is to predict either the next step or a multi-step horizon
\begin{equation}
Y_t = [x_{t+1}, \dots, x_{t+h}].
\end{equation}

\textbf{Panel forecasting.} Given observations indexed by time and region, the task is to predict a region-specific target at $t+1$ using lagged values, calendar information, and region-level attributes. This setup combines time dependence with cross-sectional heterogeneity.

These two formulations were intentionally selected to stress different inductive biases. The multivariate climate task emphasizes temporal coherence and cross-channel dependencies inside a single evolving process. The electricity task emphasizes heterogeneous entities, mixed-frequency predictors, and the need to preserve causality under realistic data availability constraints.

\subsection{Model Families}

We evaluate the following model families:
\begin{itemize}
    \item \textbf{KAN}: a two-layer KAN model using basis-weighted nonlinear transformations;
    \item \textbf{Hybrid KAN}: a linear branch plus a nonlinear KAN branch;
    \item \textbf{KANEmbedDelta}: a KAN model trained on the next-step price increment and augmented with region embeddings;
    \item \textbf{HybridKANEmbedDelta}: a strengthened hybrid KAN using region embeddings and a delta target;
    \item \textbf{ResidualRidgeKAN}: a residual-learning setup where Ridge predicts the main signal and KAN predicts the residual component;
    \item \textbf{LSTM}, \textbf{MLP}: neural baselines for climate forecasting;
    \item \textbf{Linear Regression}, \textbf{Ridge}, \textbf{HistGradientBoosting}: strong baselines for the electricity case.
\end{itemize}

For hybrid KAN models, the forecast is represented as
\begin{equation}
\hat{y} = f_{\text{linear}}(x) + f_{\text{KAN}}(x).
\end{equation}
This decomposition is useful because it allows the model to explain simple global effects through the linear branch and to capture structured nonlinearities through the KAN branch.

In the strengthened electricity setup, the residual model can also be written as
\begin{equation}
\hat{y} = f_{\text{Ridge}}(x) + f_{\text{KAN}}(x),
\end{equation}
where the second term is trained on residual targets. Conceptually, this setup shifts KAN from full-response modeling to targeted nonlinear correction.

\subsection{Implementation Notes}

The experiments were implemented in Python using NumPy, pandas, scikit-learn, PyTorch, and Jupyter notebooks for exploratory analysis. For the electricity task, the workflow was organized as a reproducible pipeline with separate stages for panel construction, leakage-safe feature alignment, model training, metric calculation, and interpretability output generation.

The KAN-based models were implemented as compact research prototypes rather than as direct wrappers around a third-party benchmark suite. This was necessary because the study required explicit control over leakage-safe temporal alignment, region embedding logic, and extraction of interpretable intermediate artifacts such as learned $\phi$-functions and grouped channel contributions.

\subsection{Interpretability Framework}

We use four complementary interpretability mechanisms:
\begin{itemize}
    \item visualization of learned $\phi$-functions in the first KAN layer;
    \item feature importance derived from KAN coefficients or activation-based contributions;
    \item grouped permutation importance for meaningful feature channels;
    \item qualitative inspection of predicted versus observed trajectories.
\end{itemize}

For pure KAN, interpretation is relatively direct because the model consists mostly of nonlinear channel-wise transformations. For hybrid KAN, interpretability becomes partial rather than absolute, since the final forecast combines linear and nonlinear components. We therefore interpret hybrid models as \emph{partially interpretable} rather than fully transparent.

\subsection{Evaluation Metrics}

We use the following metrics:
\begin{itemize}
    \item RMSE for overall predictive deviation;
    \item MAE for average absolute error;
    \item $R^2$ for explained variance;
    \item qualitative robustness comparisons under noise, missingness, and reduced sample size in the climate case.
\end{itemize}

RMSE is emphasized in the main discussion because it penalizes large errors more strongly and therefore better reflects the cost of unstable forecasts in many real settings. MAE is included because it is easier to interpret in the original target scale, while $R^2$ provides a normalized summary of how much variance is captured by the model.

\section{Case Study I: Climate Forecasting}

\subsection{Dataset}

The climate case study uses the UCI Air Quality dataset \cite{devito2008airquality}. After cleaning and removing invalid placeholders, we focus on the following five variables:
\begin{itemize}
    \item CO(GT);
    \item NO$_2$(GT);
    \item C$_6$H$_6$(GT);
    \item temperature;
    \item relative humidity.
\end{itemize}

These variables form a multivariate temporal process with both chemical and meteorological structure. The data are standardized, and sliding windows are used to construct train/test splits without shuffling, preserving temporal order.

\subsection{Preprocessing}

The climate dataset required several preprocessing steps before modeling. First, invalid placeholders representing missing sensor measurements were removed or converted to proper missing values. Second, the selected variables were aligned in time and standardized. Third, forecasting samples were created with sliding windows so that each training example contained a fixed-length history and a future prediction target.

This preprocessing is relatively conventional, but it is important for fair comparison. Without a consistent windowing strategy, the performance of sequence models such as LSTM and KAN would not be directly comparable.

\subsection{Experimental Setup}

The climate experiments include:
\begin{itemize}
    \item standard forecasting on the full dataset;
    \item forecasting with artificially injected missing values;
    \item forecasting with additive noise;
    \item forecasting with a reduced training set.
\end{itemize}

This design allows us to evaluate not only pointwise predictive quality but also robustness to data degradation, which is important for practical forecasting systems.

The climate case therefore serves two purposes in the overall paper. It is not only a benchmark where KAN can be tested under favorable sequence conditions, but also a controlled environment for analyzing the qualitative behavior of learned nonlinear transformations.

\subsection{Results on the Full Dataset}

Table~\ref{tab:climate_main} summarizes the main results on the climate forecasting task.

\begin{table}[t]
\centering
\caption{Climate forecasting results on the full dataset}
\label{tab:climate_main}
\begin{tabular}{lccc}
\toprule
Model & RMSE & MAE & $R^2$ \\
\midrule
Hybrid KAN & 0.881 & 0.696 & 0.344 \\
LSTM & 0.923 & 0.726 & 0.281 \\
MLP & 0.998 & 0.793 & 0.158 \\
KAN & 1.014 & 0.799 & 0.132 \\
\bottomrule
\end{tabular}
\end{table}

Hybrid KAN achieved the best RMSE among the tested neural architectures. This suggests that, for a multivariate sequence forecasting task with moderate dimensionality, combining a direct linear path with KAN-based nonlinear feature transformations is beneficial.

\subsection{Robustness Experiments}

Table~\ref{tab:climate_robust} reports RMSE under the robustness scenarios used in the climate experiments.

\begin{table}[t]
\centering
\caption{Climate forecasting robustness results (RMSE)}
\label{tab:climate_robust}
\begin{tabular}{lcccc}
\toprule
Model & Full & Missing & Noise & Small \\
\midrule
Hybrid KAN & 0.881 & 0.855 & 0.871 & 0.931 \\
KAN & 1.014 & 1.033 & 1.023 & 1.021 \\
LSTM & 0.923 & 0.919 & 0.943 & 0.930 \\
MLP & 0.998 & 0.948 & 0.988 & 0.943 \\
\bottomrule
\end{tabular}
\end{table}

The robustness results are important for two reasons. First, they show that the superiority of Hybrid KAN on the climate task is not limited to the clean full-data regime. Second, they support the idea that hybridization may stabilize KAN-style nonlinear transformations when the input signal is degraded.

\subsection{Interpretability on Climate Data}

The climate task is especially convenient for KAN interpretation. The learned $\phi$-functions can be visualized directly for selected input-output pairs, which helps reveal whether the model relies on monotonic, oscillatory, or saturating nonlinear transformations. In our experiments, the $\phi$-function plots showed that different output variables were associated with noticeably different input-response patterns, indicating that the network was not simply learning a generic global transformation.

From the perspective of scientific machine learning, this is one of the strongest arguments in favor of KAN: even when it is not strictly the best model in every scenario, it provides a structured lens for understanding nonlinear dependencies between input channels and predicted quantities.

\subsection{Summary of the Climate Case}

The climate experiments support three concise conclusions. First, Hybrid KAN is competitive and in our setup clearly superior to the tested LSTM and MLP baselines. Second, the KAN family remains stable under degraded data conditions. Third, the interpretability of learned $\phi$-functions gives KAN a scientific advantage over more opaque sequence models when the goal is not only accurate forecasting but also understanding variable interactions.

\section{Case Study II: Electricity Price Forecasting}

\subsection{Raw Data Sources}

The second case study focuses on forecasting the average weighted electricity purchase price in the Russian regional energy market. The dataset was constructed from three groups of sources:
\begin{itemize}
    \item hourly trading and power flow data;
    \item daily loss data;
    \item monthly tariff and cost indicators.
\end{itemize}

The hourly source contributes lagged price and demand structure, the daily source adds information about network losses, and the monthly source captures slower-moving regulatory and tariff-related effects. This combination creates a rich but challenging forecasting environment.

From the standpoint of data mining, this is a mixed-frequency fusion problem. Variables observed at different temporal resolutions must be aligned into a common prediction table without violating causal availability. This turns the data preparation stage into a substantial part of the methodological contribution.

\subsection{Leakage-Safe Dataset Construction}

A major methodological risk in energy forecasting is the unintentional inclusion of future information in present-time predictors. To avoid this, we formulated a leakage-safe prediction pipeline:
\begin{itemize}
    \item the target is defined as the purchase price at time $t+1$;
    \item daily losses are used only with a one-day lag;
    \item monthly tariff-related indicators are used only with a one-month lag.
\end{itemize}

The resulting panel covers 34 Russian regions and contains 99{,}278 observations. This setup is materially more reliable than simply joining all available features at the same timestamp, which would lead to overly optimistic evaluation.

\subsection{Feature Engineering}

The feature space included several groups of predictors:
\begin{itemize}
    \item recent lagged prices;
    \item short-term historical demand and flow-related indicators;
    \item calendar features such as hour, day-of-week, and cyclical encodings;
    \item lagged daily network loss indicators;
    \item lagged monthly tariff and cost variables.
\end{itemize}

In the strengthened setup, the feature space was deliberately compacted. The main idea was not to maximize raw feature count, but to reduce redundant predictors and make the optimization problem more favorable for KAN. In this sense, feature engineering and architecture design were tightly connected.

\subsection{Exploratory Findings}

Exploratory analysis revealed three major properties of the electricity data.

\textbf{First,} the price series displays pronounced intraday seasonality. Different hours of the day are associated with different price levels, reflecting the operational demand cycle.

\textbf{Second,} the process is highly inertial. Recent price history and recent consumption-related signals carry substantial predictive information.

\textbf{Third,} there is strong regional heterogeneity. The data do not form a single homogeneous time series; instead, they combine multiple related but distinct regional processes with different levels, variances, and temporal behavior.

These observations are crucial because they explain why a model that works well on a clean multivariate sequence may struggle in a panel setting unless the formulation is carefully adapted.

\subsection{Temporal Split Design}

The leakage-safe electricity evaluation used a chronological split. Training was performed on the early part of the panel, validation on the subsequent segment, and testing on the most recent segment. This mirrors realistic deployment conditions and avoids the false optimism that would arise from random shuffling across time.

Although rolling-origin evaluation would be even stronger methodologically, the current split already preserves the essential forecasting constraint: every prediction is generated only from information available before the target timestamp.

\subsection{Baseline Leakage-Safe Experiment}

Table~\ref{tab:elec_base} presents the results of the first leakage-safe electricity experiment.

\begin{table}[t]
\centering
\caption{Electricity forecasting: baseline leakage-safe setup}
\label{tab:elec_base}
\begin{tabular}{lccc}
\toprule
Model & RMSE & MAE & $R^2$ \\
\midrule
HistGradientBoosting & 101.32 & 55.35 & 0.944 \\
Ridge & 113.19 & 67.62 & 0.930 \\
Linear Regression & 113.29 & 67.74 & 0.930 \\
Naive Current Price & 147.70 & 85.49 & 0.881 \\
KAN & 149.51 & 97.17 & 0.878 \\
Hybrid KAN & 181.23 & 142.16 & 0.820 \\
\bottomrule
\end{tabular}
\end{table}

In this initial setup, HistGradientBoosting emerged as the strongest model. A simple interpretation is that the task is, in its raw form, closer to a strong tabular forecasting problem than to a classical multivariate neural sequence modeling problem. The poor performance of the original Hybrid KAN suggested that simply transferring KAN-style modeling from the climate setup was not sufficient.

\subsection{Strengthened KAN Formulation}

To improve KAN-based performance on electricity data, we reformulated the problem in three ways.

\textbf{1) Delta target.} Instead of predicting the absolute next-step price, the model predicts
\begin{equation}
\Delta p_t = p_{t+1} - p_t.
\end{equation}
This transformation reduces level effects and makes the problem closer to local dynamics modeling.

\textbf{2) Region embedding.} Instead of only treating the region as a sparse categorical variable, we represent each region with a trainable embedding vector. This makes it easier for the model to learn region-specific offsets and latent structure.

\textbf{3) Residual learning.} We use a Ridge baseline to explain the main linear component and then train KAN to model the remaining residual signal. This is particularly appealing because it allows KAN to focus on nonlinear effects rather than on the full target from scratch.

\subsection{Results After Strengthening}

The strengthened setup substantially improved KAN-family performance, as shown in Table~\ref{tab:elec_strong}.

\begin{table}[t]
\centering
\caption{Electricity forecasting: strengthened KAN setup}
\label{tab:elec_strong}
\begin{tabular}{lccc}
\toprule
Model & RMSE & MAE & $R^2$ \\
\midrule
HGBDelta & 102.56 & 55.27 & 0.943 \\
HybridKANEmbedDelta & 107.57 & 61.15 & 0.937 \\
ResidualRidgeKAN & 108.56 & 60.81 & 0.936 \\
KANEmbedDelta & 115.17 & 69.88 & 0.928 \\
RidgeDelta & 116.30 & 69.39 & 0.926 \\
\bottomrule
\end{tabular}
\end{table}

The strongest KAN-based model became HybridKANEmbedDelta with RMSE 107.57, reducing the gap to the best boosting baseline to approximately 5 RMSE points. This is a dramatic improvement compared to the original Hybrid KAN result of RMSE 181.23.

\subsection{Relative Improvement}

To emphasize the scale of the improvement, Table~\ref{tab:improvement} compares the original and strengthened KAN-family models.

\begin{table}[t]
\centering
\caption{Improvement of KAN-family models on electricity data}
\label{tab:improvement}
\begin{tabular}{lcc}
\toprule
Model family & Original RMSE & Strengthened RMSE \\
\midrule
KAN $\rightarrow$ KANEmbedDelta & 149.51 & 115.17 \\
Hybrid KAN $\rightarrow$ HybridKANEmbedDelta & 181.23 & 107.57 \\
\bottomrule
\end{tabular}
\end{table}

This result is important for the overall narrative of the paper. It shows that KAN's practical value depends strongly on problem formulation. A weak result in a raw setup does not necessarily imply that the architecture itself is unsuitable; rather, it may indicate that the model needs a formulation aligned with its strengths.

\subsection{Interpretability on Electricity Data}

Interpretability remains a major reason to study KAN in the electricity setting. In our experiments, the strongest KAN-related importance patterns consistently involved:
\begin{itemize}
    \item recent price history;
    \item calendar structure such as hour-of-day and cyclical encodings;
    \item load and demand-related history.
\end{itemize}

This is an intuitively plausible result and increases confidence that the model is learning meaningful dynamics rather than relying on purely spurious relationships.

At the same time, it is important to state that interpretability differs between model types. Pure KAN models remain the most directly interpretable. Hybrid KAN models are still useful for interpretation, but only partially so, since the final prediction combines an explicit linear component and a nonlinear KAN component.

Nevertheless, even partial interpretability is valuable. In the strengthened hybrid setting, we can still inspect which inputs dominate the nonlinear correction branch, which channels contribute most to final prediction quality, and whether the learned $\phi$-functions behave in a smooth and plausible way. This is substantially more informative than treating the model as a fully opaque predictor.

\section{Discussion}

\subsection{What We Learned from the Two Cases}

The climate and electricity cases lead to complementary conclusions.

For the climate dataset, Hybrid KAN behaves as a genuinely strong forecasting architecture. The task structure is favorable for KAN: the data form a coherent multivariate sequence, the variables are closely related, and the temporal window formulation aligns well with channel-wise nonlinear transformations.

For the electricity dataset, the situation is different. The task is not simply multivariate sequence prediction; it is panel forecasting with strong cross-sectional heterogeneity and substantial dependence on engineered lag and calendar features. In this regime, tabular baselines such as gradient boosting are naturally very strong. Therefore, it is not surprising that the best overall model is not KAN.

Taken together, the two case studies argue against simplistic claims of the form ``KAN is always better'' or ``KAN is not useful in practice.'' Neither statement is supported by the evidence. The more defensible conclusion is that KAN is highly sensitive to task structure, target design, and the relationship between temporal dynamics and cross-sectional heterogeneity.

\subsection{Why KAN Still Matters}

Even though KAN is not the strongest model overall in the electricity case, it remains valuable for several reasons:
\begin{itemize}
    \item it provides structured nonlinear modeling;
    \item it offers more interpretable internal transformations than typical black-box deep models;
    \item with appropriate target design and feature representation, it becomes much more competitive;
    \item it can be integrated into hybrid and residual pipelines rather than used only as a standalone predictor.
\end{itemize}

Thus, KAN should be understood less as a universal winner and more as a model family that is especially useful when interpretability is important and when the problem can be formulated in a way that matches its inductive bias.

\subsection{Practical Implications}

For practical forecasting projects, the results suggest a useful workflow:
\begin{itemize}
    \item start with strong linear and boosting baselines;
    \item introduce KAN when nonlinear interpretability is important;
    \item prefer hybrid or residual KAN designs over naive direct transfer from unrelated tasks;
    \item pay particular attention to target definition and feature grouping.
\end{itemize}

This workflow is valuable for thesis work and applied research because it frames KAN not as a replacement for all other methods, but as part of a broader model design toolkit.

\subsection{Limits of the Current Study}

This work has several limitations.

First, the climate experiments are comparatively smaller and cleaner than the electricity case, which makes direct cross-domain conclusions impossible. Second, the electricity data span a relatively limited time range after applying leakage-safe intersection rules. Third, although the strengthened KAN setup significantly improved results, it did not surpass the best boosting baseline.

Moreover, the study focuses on a specific family of KAN implementations and does not exhaustively cover all possible sequence-oriented KAN designs, attention-based hybrids, or transformer-style variants.

\section{Threats to Validity}

\textbf{Internal validity.} For the electricity case, the main threat was leakage from future information. We explicitly addressed this through lagged daily and monthly features and by forecasting $t+1$ rather than using contemporaneous signals that could encode future knowledge.

\textbf{Construct validity.} RMSE, MAE, and $R^2$ capture predictive quality but do not fully capture operational utility. In some applications, ranking errors, directional errors, or tail-risk performance might matter more.

\textbf{External validity.} The climate and electricity datasets are representative of two practically meaningful settings, but they do not cover all forecasting scenarios. Therefore, the conclusions should be interpreted as evidence about \emph{patterns of applicability} rather than universal claims.

\textbf{Statistical conclusion validity.} The reported results rely on a limited number of benchmark datasets and finite hyperparameter exploration. It is therefore possible that stronger tuning of either KAN or competing baselines would shift some absolute numbers. However, the relative pattern of findings is stable enough to support the main qualitative conclusions.

\section{Conclusion}

This paper presented an extended empirical study of KAN and Hybrid KAN architectures on two forecasting tasks with substantially different structure. On the climate forecasting task, Hybrid KAN achieved the best quality among the tested neural models and demonstrated strong robustness to noisy and incomplete data. On the electricity panel forecasting task, tree-based baselines remained the strongest overall predictors. However, after strengthening the KAN formulation using delta-target prediction, region embeddings, and residual learning, the best KAN-based model approached the strongest baseline much more closely.

The central conclusion is that KAN is most valuable not as a universal forecasting winner, but as an interpretable nonlinear model family whose strengths become particularly visible under appropriate task design. For researchers and practitioners, this means that KAN should be considered both as a predictive model and as an analysis tool for studying nonlinear feature-channel effects in forecasting problems.

Future work includes region-level multi-step electricity forecasting, comparison with standard long-horizon benchmark datasets such as ETT, deeper ablation of basis size and hidden dimensionality, and exploring residual KAN designs on top of strong tree-based baselines.

\section*{Reproducibility Note}

The study was conducted in a reproducible computational workflow based on Python scripts and Jupyter notebooks. The project includes separate notebooks for exploratory analysis and separate scripts for the baseline leakage-safe electricity experiment and the strengthened KAN experiments. Output artifacts include metric tables, prediction files, and interpretability plots. This structure makes it possible to reuse the study both as a research draft and as a basis for thesis presentation materials.

\bibliographystyle{IEEEtran}
\bibliography{references}

\end{document}
